\section{结论}
\label{sec:conclusion}

\method{} 为固定预测器提供了一个推理时上下文集成接口。共识加权聚合彼此支持的预测，任务特定门控则使更新幅度与查询级证据相匹配。受控扰动、分布偏移、合成机制、自然偏移以及常规预测器上的实验表明，该接口能够在无需重新训练主干模型的情况下提高预测稳定性。

\paragraph{局限性与未来工作。}
Shared 不需要变化类别标识，而先验知情的 Specific 则研究另一种互补设置，即在推理前已知变化类别。当前实验覆盖了受控变化与若干自然偏移；更广泛的自然场景与高风险应用仍有待评估。评估多个上下文会增加推理成本，因此未来可研究根据无标签目标数据选择上下文，并进一步减少重复计算。

\section*{可复现性声明}

可复现性信息分布于正文与附录中。附录~\ref{app:theory} 给出完整理论分析；附录~\ref{app:experimental_setup} 说明数据集、划分、扰动、分布偏移、上下文构造、聚合方式与实现设置；附录~\ref{app:backbone-strength-curves} 和~\ref{app:supplemental_ablations} 分别报告按强度展开的完整结果与消融实验。所有实验结果均遵循附录~\ref{app:experimental_setup} 中的固定方案。受许可证限制的模型权重均来自其公开发布版本，本文不对其进行再分发。

\section*{人工智能工具使用声明}

作者使用生成式人工智能工具辅助润色写作、起草部分稿件，并支持投稿表格中未单独列出的任务。所有由人工智能辅助生成的内容均经作者审查与核验，作者对最终稿件承担责任。
